\documentclass[a4paper,12pt]{ctexart}
\usepackage{geometry}
\usepackage{titlesec}
\usepackage{fancyhdr}
\usepackage{amsmath}
\usepackage{graphicx}
\usepackage{listings}
\usepackage{xcolor}
\usepackage{booktabs}
\usepackage{array}
\usepackage{float}
\usepackage{enumitem}
\usepackage{ulem} % 下划线支持
\usepackage{caption}

% 页眉页脚设置
\pagestyle{fancy}
\fancyhf{}
\fancyfoot[C]{\thepage}
\renewcommand{\headrulewidth}{0pt}

% 页面设置
\geometry{left=2.5cm, right=2.5cm, top=2.5cm, bottom=2.5cm}

% 代码块设置
\definecolor{codegreen}{rgb}{0,0.6,0}
\definecolor{codegray}{rgb}{0.5,0.5,0.5}
\definecolor{codepurple}{rgb}{0.58,0,0.82}
\definecolor{backcolour}{rgb}{0.95,0.95,0.92}

\lstdefinestyle{mystyle}{
    backgroundcolor=\color{backcolour},   
    commentstyle=\color{codegreen},
    keywordstyle=\color{magenta},
    numberstyle=\tiny\color{codegray},
    stringstyle=\color{codepurple},
    basicstyle=\ttfamily\footnotesize,
    breakatwhitespace=false,         
    breaklines=true,                 
    captionpos=b,                    
    keepspaces=true,                 
    numbers=left,                    
    numbersep=5pt,                  
    showspaces=false,                
    showstringspaces=false,
    showtabs=false,                  
    tabsize=4,
    language=C++
}
\lstset{style=mystyle}

% 标题设置
\setcounter{secnumdepth}{-2} % 禁止所有自动编号，因为已使用手动编号
\titleformat{\section}{\large\bfseries\raggedright}{}{0em}{}
\titleformat{\subsection}{\normalsize\bfseries\raggedright}{}{0em}{}
\titleformat{\subsubsection}{\small\bfseries\raggedright}{}{0em}{}
\titlespacing*{\section}{0pt}{0pt}{10pt}
\titlespacing*{\subsection}{0pt}{5pt}{5pt}
\titlespacing*{\subsubsection}{0pt}{5pt}{5pt}

\begin{document}

% 封面/头部区域
\begin{center}
    \Large \textbf{中国计量大学信息工程学院人工智能系} \\
    \vspace{0.5cm}
    \textbf{实验报告}
\end{center}

\vspace{1cm}

\noindent
\begin{tabular}{p{0.12\textwidth}p{0.35\textwidth}p{0.12\textwidth}p{0.3\textwidth}}
    \textbf{实验课程：} & \underline{\makebox[0.3\textwidth][c]{算法与数据结构}} & \textbf{实验名称：} & \underline{\makebox[0.25\textwidth][c]{exp4-NMS算法}} \\
    & & & \\
    \textbf{班\quad 级：} & \underline{\makebox[0.3\textwidth][c]{24智能2班}} & \textbf{学\quad 号：} & \underline{\makebox[0.25\textwidth][c]{2400305209}} \\
    & & & \\
    \textbf{姓\quad 名：} & \underline{\makebox[0.3\textwidth][c]{史航安}} & \textbf{实验日期：} & \underline{\makebox[0.25\textwidth][c]{2025/12/11}} \\
\end{tabular}

\vspace{1cm}

\section*{一、实验目的}

1. \textbf{深入理解排序算法}：通过亲手实现快速排序、归并排序、堆排序、插入排序和冒泡排序，掌握不同排序算法的核心思想、时间复杂度特征及其在不同规模数据下的表现。
2. \textbf{掌握NMS算法原理与应用}：学习非极大值抑制（Non-Maximum Suppression, NMS）算法在计算机视觉目标检测中的作用，理解其基于置信度排序和IoU（交并比）计算的筛选机制。
3. \textbf{性能分析与优化}：探究不同排序算法对NMS整体性能的影响，分析数据规模（$N$）与数据分布（随机 vs 聚集）对算法效率的影响；设计并实现基于栅格分桶的优化版NMS，对比其与基础版NMS的性能差异。
4. \textbf{工程实践能力}：提升使用C++及STL（Standard Template Library）进行算法设计、模块化编程、性能测试及构建工具（CMake）使用的综合能力。

\section*{二、实验内容}

1. \textbf{排序算法实现}：
   \begin{itemize}
       \item 基于自定义 \texttt{Vector} 容器实现五种排序算法：快速排序（Quick Sort）、归并排序（Merge Sort）、堆排序（Heap Sort）、插入排序（Insertion Sort）、冒泡排序（Bubble Sort）。
       \item 实现比较器（Comparator），支持按边界框（Box）的置信度（Score）进行降序排列。
   \end{itemize}

2. \textbf{NMS算法实现}：
   \begin{itemize}
       \item \textbf{基础NMS}：实现标准的 $O(N^2)$ NMS流程，包括排序、IoU计算与重叠抑制。
       \item \textbf{栅格NMS}：设计基于空间网格索引（Grid-based）的优化NMS，通过将边界框映射到二维网格减少无效的IoU计算，提升大规模数据下的处理效率。
   \end{itemize}

3. \textbf{数据生成}：
   \begin{itemize}
       \item 设计 \textbf{随机分布} 生成器：在画布上均匀随机生成边界框。
       \item 设计 \textbf{聚集分布} 生成器：基于多个中心点生成高斯分布的簇状边界框，模拟真实场景中目标密集的分布。
       \item 支持生成不同规模的数据集（$N=100, 1000, 3000$）。
   \end{itemize}

4. \textbf{性能测试}：
   \begin{itemize}
       \item 记录每种排序算法在不同NMS实现（基础 vs 栅格）、不同数据分布、不同数据规模下的运行时间。
       \item 输出格式化的性能对比报告。
   \end{itemize}

\section*{三、实验环境}

\begin{itemize}
    \item \textbf{操作系统}：Windows 11
    \item \textbf{开发工具}：Trae IDE / VS Code
    \item \textbf{编译器}：MinGW GCC / MSVC (支持 C++17)
    \item \textbf{构建系统}：CMake 3.16+
    \item \textbf{硬件环境}：Intel Core i7 / 16GB RAM (示例配置)
    \item \textbf{核心库}：MySTL (本地自研简易STL库，包含Vector等实现)
\end{itemize}

\vspace{2cm}

\noindent
\begin{tabular}{p{0.45\textwidth}p{0.45\textwidth}}
    \textbf{实验成绩：} \underline{\hspace{4cm}} & \textbf{主管教师签名：} \underline{\hspace{4cm}} \\
\end{tabular}

\newpage

\section*{四、实验数据记录（源程序或算法设计思想）}

\subsection{1. 数据结构定义}

实验核心数据结构为 \texttt{Box}，包含左上角坐标 $(x_1, y_1)$、右下角坐标 $(x_2, y_2)$、置信度 $score$ 以及原始索引 $id$。

\subsubsection{IoU 计算实现}
\begin{lstlisting}
struct Box {
    float x1, y1, x2, y2, score;
    int id;
};

// IoU calculation function
static inline float iou(const Box& a, const Box& b) {
    // Calculate intersection area
    float ix1 = std::max(a.x1, b.x1);
    float iy1 = std::max(a.y1, b.y1);
    float ix2 = std::min(a.x2, b.x2);
    float iy2 = std::min(a.y2, b.y2);
    float iw = ix2 - ix1;
    float ih = iy2 - iy1;
    if (iw <= 0 || ih <= 0) return 0.0f; // No intersection
    
    float inter = iw * ih;
    float areaA = (a.x2 - a.x1) * (a.y2 - a.y1);
    float areaB = (b.x2 - b.x1) * (b.y2 - b.y1);
    float uni = areaA + areaB - inter; // Union area
    
    if (uni <= 0) return 0.0f;
    return inter / uni;
}
\end{lstlisting}
\subsection{2. 排序算法设计}

本次实验实现了五种排序算法，通过模板函数与自定义比较器实现对 \texttt{Box} 对象的通用排序。

\subsubsection{(1) 快速排序 (Quick Sort)}
采用分治策略，选择 pivot（本实现选取区间末尾元素），将数组划分为两部分，左侧元素均优于 pivot，右侧元素均劣于 pivot，递归处理。
\begin{lstlisting}
template<typename T, typename Compare>
static int part(Vector<T>& a, int lo, int hi, Compare comp) {
    T pivot = a[hi - 1];
    int i = lo;
    for (int j = lo; j < hi - 1; ++j) {
        if (comp(pivot, a[j])) { // Descending comparison: a[j] > pivot
            vswap(a[i], a[j]); ++i; 
        }
    }
    vswap(a[i], a[hi - 1]);
    return i;
}
\end{lstlisting}

\subsubsection{(2) 冒泡排序 (Bubble Sort)}
重复遍历列表，比较相邻元素，若顺序错误则交换，直到没有交换发生。本实验中为了观察 $O(N^2)$ 排序对 NMS 的影响而加入。
\begin{lstlisting}
template<typename T, typename Compare>
static void sort_bubble(Vector<T>& a, Compare comp) {
    int n = a.size();
    bool swapped = true;
    while (swapped) {
        swapped = false;
        for (int i = 1; i < n; ++i) {
            if (comp(a[i - 1], a[i])) { // If prev < next (lower score), swap
                vswap(a[i - 1], a[i]); 
                swapped = true; 
            }
        }
        --n; 
        if (n <= 1) break;
    }
}
\end{lstlisting}

\subsubsection{(3) 堆排序 (Heap Sort)}
利用堆数据结构（本例隐含构建大顶堆），将堆顶元素（最大值）与末尾交换，并缩小堆范围进行下滤（Heapify）。由于需要降序，实际操作中逻辑稍作调整或最终反转。

\subsubsection{(4) 归并排序 (Merge Sort)}
采用分治法，将序列二分递归排序，再合并两个有序子序列。稳定性好，但需要 $O(N)$ 额外空间。

\subsubsection{(5) 插入排序 (Insertion Sort)}
构建有序序列，对于未排序数据，在已排序序列中从后向前扫描，找到相应位置并插入。小规模数据表现尚可，大规模下性能较差。

\subsection{3. NMS 算法设计}

\subsubsection{(1) 基础 NMS (Basic NMS)}
算法流程：
\begin{enumerate}
    \item 按置信度对所有框进行排序（降序）。
    \item 初始化所有框为“未抑制”状态。
    \item 遍历排序后的列表：
    \begin{itemize}
        \item 若当前框 $A$ 已被抑制，跳过。
        \item 保留框 $A$。
        \item 遍历 $A$ 之后的所有框 $B$：
        \begin{itemize}
            \item 若 $B$ 未被抑制且 $\text{IoU}(A, B) > \text{threshold}$，则抑制 $B$。
        \end{itemize}
    \end{itemize}
\end{enumerate}
时间复杂度：排序 $O(N \log N)$ + 抑制 $O(N^2)$，总体 $O(N^2)$。

\subsubsection{(2) 栅格分桶 NMS (Grid NMS)}
算法流程：
\begin{enumerate}
    \item 按置信度排序。
    \item 将图像划分为 $W \times H$ 的网格（例如 $50 \times 50$ 像素为一格）。
    \item 预处理：将每个框根据中心点坐标放入对应的网格桶中。
    \item 遍历排序后的列表：
    \begin{itemize}
        \item 若当前框 $A$ 已被抑制，跳过。
        \item 保留框 $A$。
        \item 计算 $A$ 中心所在的网格及其 3x3 邻域网格。
        \item \textbf{仅} 遍历邻域网格中的框 $B$：
        \begin{itemize}
            \item 若 $B$ 的索引在 $A$ 之后且未被抑制，计算 IoU 并进行抑制。
        \end{itemize}
    \end{itemize}
\end{enumerate}
优势：将抑制阶段的全局搜索缩小为局部搜索，大幅降低比较次数。

\begin{lstlisting}
static Vector<Box> nms_grid(Vector<Box> boxes, float thr, SortAlgo algo, float cell) {
    apply_sort(boxes, algo);
    // ... Calculate grid dimensions and allocate ...
    
    // Build grid index
    Vector<Vector<int>> grid;
    // ... Insert each box index into corresponding grid[idx] ...

    Vector<Box> keep;
    // ... Traverse each box ...
        // Search potential overlapping boxes ONLY in 3x3 neighbor grids
        for (int yy = gy0; yy <= gy1; ++yy) {
            for (int xx = gx0; xx <= gx1; ++xx) {
                int idx = yy * gw + xx;
                for (int t = 0; t < grid[idx].size(); ++t) {
                    // ... Compare IoU and suppress ...
                }
            }
        }
    return keep;
}
\end{lstlisting}

\section*{五、实验数据分析及结论（程序运行结果及分析）}

\subsection{1. 实验结果数据}

以下数据基于 $1000 \times 1000$ 画布，IoU 阈值 0.5，栅格大小 50 的测试环境。单位：毫秒 (ms)。

\begin{table}[H]
\centering
\caption{不同算法在随机分布下的NMS耗时对比}
\label{tab:random_dist}
\begin{tabular}{@{}lcccc@{}}
\toprule
\textbf{算法类型} & \textbf{规模 N=100} & \textbf{规模 N=1000} & \textbf{规模 N=3000} & \textbf{时间复杂度} \\ \midrule
\textbf{基础 NMS} & & & & \\
快速排序 & 0 & 9 & 66 & $O(N^2)$ \\
归并排序 & 0 & 9 & 66 & $O(N^2)$ \\
堆排序   & 0 & 10 & 69 & $O(N^2)$ \\
插入排序 & 0 & 9 & 87 & $O(N^2)$ \\
冒泡排序 & 0 & 13 & 112 & $O(N^2)$ \\ \midrule
\textbf{栅格 NMS} & & & & \\
快速排序 & 0 & 1 & 4 & $O(N \log N + K \cdot N)$ \\
归并排序 & 1 & 1 & 5 & $\approx O(N)$ (近似) \\
堆排序   & 0 & 1 & 5 & \\
插入排序 & 1 & 3 & 22 & \\
冒泡排序 & 1 & 6 & 43 & \\ \bottomrule
\end{tabular}
\end{table}

\begin{table}[H]
\centering
\caption{不同算法在聚集分布下的NMS耗时对比}
\label{tab:cluster_dist}
\begin{tabular}{@{}lcccc@{}}
\toprule
\textbf{算法类型} & \textbf{规模 N=100} & \textbf{规模 N=1000} & \textbf{规模 N=3000} & \textbf{性能分析} \\ \midrule
\textbf{基础 NMS} & & & & \\
快速排序 & 0 & 7 & 54 & \\
插入排序 & 0 & 11 & 68 & \\
冒泡排序 & 0 & 13 & 98 & \\ \midrule
\textbf{栅格 NMS} & & & & \\
快速排序 & 0 & 2 & 13 & 局部搜索优势 \\
插入排序 & 0 & 4 & 31 & \\
冒泡排序 & 0 & 9 & 57 & \\ \bottomrule
\end{tabular}
\end{table}

\subsection{2. 结果分析}

\subsubsection{(1) 排序算法的影响}
在基础 NMS 中，整体耗时主要由 $O(N^2)$ 的抑制过程主导。
\begin{itemize}
    \item 当 $N=3000$ 时，快速、归并、堆排序差异不大（约 66ms），因为排序本身的 $O(N \log N)$ 耗时占比相对较小。
    \item 插入排序（77ms）和冒泡排序（112ms）明显更慢，尤其是冒泡排序，其 $O(N^2)$ 的排序开销与 NMS 的抑制开销叠加，导致整体性能最差。
    \item \textbf{结论}：对于 NMS，选择 $O(N \log N)$ 的排序算法（如快速排序）是必要的，特别是当 $N$ 较大时。
\end{itemize}

\subsubsection{(2) 基础 NMS vs 栅格 NMS}
\begin{itemize}
    \item \textbf{性能提升显著}：栅格 NMS 在 $N=3000$ 时将耗时从约 66ms 降低到了 4-5ms，加速比超过 10 倍。这是因为基础 NMS 需要计算几乎所有未抑制框对的 IoU，而栅格 NMS 仅计算空间邻域内的框。
    \item \textbf{复杂度优化}：栅格 NMS 成功将抑制阶段的复杂度从 $O(N^2)$ 降低到了接近线性的水平（取决于数据密度和栅格大小）。
    \item \textbf{排序瓶颈凸显}：在栅格 NMS 中，由于抑制阶段极快，排序阶段的耗时占比大幅提升。例如，冒泡排序在栅格 NMS 下的总耗时（43ms）远高于快速排序（4ms），这清晰地展示了低效排序算法对优化后系统的拖累。
\end{itemize}

\subsubsection{(3) 数据分布的影响}
\begin{itemize}
    \item \textbf{聚集分布}：在聚集分布下，框的重叠率更高。基础 NMS 在聚集分布下（54ms）略快于随机分布（66ms），这可能是因为高分框能更快地抑制大量重叠框，从而减少了后续遍历的次数（被抑制的框不需要再去抑制别人）。
    \item \textbf{栅格效率}：聚集分布下栅格 NMS 的耗时（13ms）略高于随机分布（4ms），这是因为聚集区域内的网格桶中包含更多元素，导致局部邻域内的比较次数增加。
\end{itemize}

\subsection{3. 结论}
\begin{enumerate}
    \item \textbf{算法选择}：实际工程中应首选 \textbf{快速排序 + 栅格/索引优化 NMS} 的组合。
    \item \textbf{优化策略}：当数据规模增长时，单纯优化排序算法收益有限（受限于 $O(N^2)$ 的抑制瓶颈）；引入空间索引（如栅格）打破 $O(N^2)$ 瓶颈是性能提升的关键。
    \item \textbf{分布敏感性}：算法性能受数据分布影响，聚集型数据虽然利于基础 NMS 的“早停”，但会增加栅格 NMS 的桶内负载，需根据场景调整栅格大小。
\end{enumerate}

\section*{六、实验中的问题及心得体会}

\subsection{1. 遇到的问题}
1. \textbf{中文乱码问题}：在 Windows 控制台输出实验结果时，初期出现乱码。
   - \textbf{解决}：引入 \texttt{windows.h} 并在 \texttt{main} 函数中设置 \texttt{SetConsoleOutputCP(CP\_UTF8)}，强制控制台使用 UTF-8 代码页。
2. \textbf{栈溢出风险}：在实现快速排序和归并排序的递归版本时，考虑到若数据量极大（如 $N > 10^5$），可能导致栈溢出。
   - \textbf{思考}：虽然本次实验 $N \le 3000$ 未出现此问题，但生产环境应考虑改用非递归实现或 introspective sort。
3. \textbf{栅格边界处理}：在实现栅格 NMS 时，处理边界框中心点坐标映射到网格索引时，需小心处理数组越界（例如坐标正好等于图像宽高的情况）。
   - \textbf{解决}：增加了 \texttt{clamp} 逻辑，确保 \texttt{gx, gy} 始终在 \texttt{[0, gw-1]} 范围内。

\subsection{2. 心得体会}
通过本次实验，我不仅复习了经典的排序算法，更深刻体会到了**“算法复杂度”**在实际运行时间上的具体体现。
- **理论与实践结合**：看着 $O(N^2)$ 的冒泡排序在数据量增加时耗时剧增，而 $O(N \log N)$ 算法保持稳健，验证了理论分析的正确性。
- **优化的层次**：优化不能只盯着某一个环节。在 NMS 中，优化排序只能提升一部分性能，而改变算法逻辑（引入栅格索引）降低整体复杂度阶数，才能带来质的飞跃。这启示我在未来的开发中，要先分析系统的\textbf{主要瓶颈}，再对症下药。
- **工程细节**：从 CMake 构建脚本的编写到 Windows 字符集的处理，这些工程细节虽然琐碎，但却是保证程序正确运行和跨平台移植的基础。

\end{document}
